\section*{Chapter \ref{ch:b2:crustal-melting} --- Crustal Melting: Migmatites and Granites}

\begin{solution}{exo:b2:crustal-melting:1}
(a) Per \qty{100}{\gram}: $n(\ce{Al2O3}) = 17.03/101.96 = \qty{0.16703}{mol}$,
$n(\ce{CaO}) = 0.09183$, $n(\ce{Na2O}) = 0.06783$, $n(\ce{K2O}) = 0.03076$.
A/CNK $= 0.16703/0.19042 = 0.877$; A/NK $= 0.16703/0.09859 = 1.694$.
(b) A/CNK $< 1 <$ A/NK: \omterm{def:b2:crustal-melting:asi}{metaluminous}. The calcium the feldspars cannot take goes into
calcium silicates: hornblende (or clinopyroxene, titanite).
\end{solution}

\begin{solution}{exo:b2:crustal-melting:2}
(a) Muscovite and garnet hold aluminium beyond the feldspars: \omterm{def:b2:crustal-melting:asi}{peraluminous}; this granite
can be an S-type. (b) Hornblende and titanite hold calcium beyond the feldspars:
\omterm{def:b2:crustal-melting:asi}{metaluminous}. (c) Aegirine and a sodic amphibole hold sodium without the aluminium
of feldspar: \omterm{def:b2:crustal-melting:asi}{peralkaline}.
\end{solution}

\begin{solution}{exo:b2:crustal-melting:3}
One kilometre of rock weighs $2800\times9.81\times1000 = \qty{2.747e7}{Pa}$, so
$P = 0.2747\,z$ (\unit{kbar}, $z$ in \unit{km}). The solidus line is $T = 680 - 7.5(P -
2) = 695 - 2.060\,z$. Setting $10 + 20z = 695 - 2.060z$ gives $z = 685/22.060 =
\qty{31.1}{km}$, $T = \qty{631}{\celsius}$, $P = \qty{8.5}{kbar}$.
\end{solution}

\begin{solution}{exo:b2:crustal-melting:4}
(a) The white layers are \omterm{def:b2:crustal-melting:migmatite}{leucosomes}, the black seams \omterm{def:b2:crustal-melting:migmatite}{melanosomes}, the ordinary gneiss the
mesosome. (b) Discrete, well-defined \omterm{def:b2:crustal-melting:migmatite}{leucosomes}, \omterm{def:b2:crustal-melting:migmatite}{melanosomes} and mesosomes: a \omterm{def:b2:crustal-melting:migmatite}{metatexite},
with a stromatic structure. (c) The rock has become a \omterm{def:b2:crustal-melting:migmatite}{diatexite} (schlieric and nebulitic):
melting went further, the dark minerals melted too, and the melt stayed in place,
erasing the layering.
\end{solution}

\begin{solution}{exo:b2:crustal-melting:5}
(a) $M = 39.10 + 3\times24.31 + 26.98 + 3\times28.09 + 12\times16.00 + 2\times1.008 =
\qty{417.30}{\gram\per\mole}$; $w = 18.02/417.30 = \qty{4.32}{\percent}$.
(b) $F = 0.25\times0.0432/0.040 = 0.270$, or \qty{27}{\percent} of melt by mass: within the
20 to \qty{30}{\percent} by volume produced by biotite dehydration-melting experiments on
similar rocks.
\end{solution}

\begin{solution}{exo:b2:crustal-melting:6}
(a) K-feldspar \qty{278.35}{\gram\per\mole}, albite 262.24, anorthite 278.22, muscovite
398.33. (b) Amounts: K-feldspar $30/278.35 = \qty{0.10778}{mol}$, albite $0.13347$,
anorthite $0.01797$, muscovite $0.01255$. K-feldspar is $\tfrac12$\ce{K2O.}$\tfrac12$\ce{Al2O3},
albite $\tfrac12$\ce{Na2O.}$\tfrac12$\ce{Al2O3}, anorthite \ce{CaO.Al2O3}, muscovite
$\tfrac12$\ce{K2O.}$\tfrac32$\ce{Al2O3} per formula unit. $n(\ce{Al2O3}) = 0.05389 +
0.06673 + 0.01797 + 0.01883 = \qty{0.15742}{mol}$; $n(\ce{K2O}) = 0.05389 + 0.00628 =
0.06017$, $n(\ce{Na2O}) = 0.06673$, $n(\ce{CaO}) = 0.01797$, sum $0.14487$.
A/CNK $= 0.15742/0.14487 = 1.087$. (c) Without the muscovite the three feldspars give
exactly 1 (\cref{prop:b2:crustal-melting:feldspar}); the muscovite, with one \ce{Al2O3}
more than K-feldspar per formula unit, makes the rock \omterm{def:b2:crustal-melting:asi}{peraluminous}.
\end{solution}

\begin{solution}{exo:b2:crustal-melting:7}
Basalt: $900 + 2.5\times(1200 - 850) = \qty{1775}{\mega\joule\per\cubic\metre}$. Crust:
$2.5\times(850 - 650) + 0.20\times900 = \qty{680}{\mega\joule\per\cubic\metre}$.
$V_c/V_b = 1775/680 = 2.6$. The estimate ignores the heat that leaks away by conduction
upwards and sideways and the time heat takes to spread; it also assumes that all the
basalt's heat reaches crust that is ready to melt: it is an upper bound.
\end{solution}

\begin{solution}{exo:b2:crustal-melting:8}
(a) $n(\ce{Al2O3}) = 14.15/101.96 = 0.13878$, $n(\ce{CaO}) = 0.01753$, $n(\ce{Na2O}) =
0.07809$, $n(\ce{K2O}) = 0.04416$; A/CNK $= 0.13878/0.13978 = 0.993$, A/NK $=
0.13878/0.12225 = 1.135$: \omterm{def:b2:crustal-melting:asi}{metaluminous}, just. (b) With \qty{14.30}{\percent} \ce{Al2O3},
$0.14025/0.13978 = 1.003$; with \qty{14.00}{\percent}, $0.13731/0.13978 = 0.982$. The index
straddles 1 within the certified uncertainty: the obsidian lies on the
metaluminous--peraluminous boundary, and its class is not robust.
\end{solution}

\begin{solution}{exo:b2:crustal-melting:9}
(a) Means $\bar P = \qty{5.0}{kbar}$, $\bar w = \qty{10.045}{\percent}$;
$\sum(P - \bar P)(w - \bar w) = 17.62$, $\sum(P - \bar P)^2 = 14.0$; slope
$17.62/14.0 = \qty{1.259}{\percent}$ per kilobar, intercept $10.045 - 1.259\times5.0 =
\qty{3.75}{\percent}$. The slope agrees with the quoted 1.27 to within
\qty{1}{\percent}. (b) $3.75 + 1.259\times6 = \qty{11.30}{\percent}$. (c) $10/0.1130 =
\qty{88.5}{\gram}$ of melt. (d) The line rests on data between 3 and \qty{8}{kbar} at
\qty{800}{\celsius} only; nothing guarantees that solubility stays linear at higher
pressure, where the melt and the fluid become more alike.
\end{solution}

\begin{solution}{exo:b2:crustal-melting:10}
(a) The line rises by \qty{4}{kbar} over \qty{50}{\celsius}: $P = 6 + 0.08\,(T - 750)$;
at \qty{770}{\celsius}, $P = \qty{7.6}{kbar}$. (b) From 7.6 to \qty{6}{kbar} the
temperature falls along the line from 770 to \qty{750}{\celsius}; the heat released,
$2.5\times20 = \qty{50}{\mega\joule\per\cubic\metre}$, melts $50/900 = 0.056$, about
\qty{6}{\percent} of the rock by volume. (c) This is the window proposed for the
Himalayan \omterm{def:b2:crustal-melting:leucogranite}{leucogranites}, 6 to \qty{8}{kbar} and 750 to \qty{770}{\celsius}: decompression
alone makes melt, but only a few per cent per crossing of the curve.
\end{solution}

\begin{solution}{exo:b2:crustal-melting:11}
$\ell^2/\kappa = (2000)^2/10^{-6} = \qty{4e12}{s} = \qty{0.13}{Myr}$ for the \omterm{def:b2:magma-emplacement:bodies}{sill};
$(10^4)^2/10^{-6} = \qty{1e14}{s} = \qty{3.2}{Myr}$ for the underplate. Both are short
beside ten million years: one event cannot keep the crust molten that long, which needs
heat supplied again and again, by repeated intrusions or by the slow heating of a
thickened crust.
\end{solution}

\begin{solution}{exo:b2:crustal-melting:12}
(a) $5.0x + 2.0(1 - x) = 2.75$ gives $x = 0.75/3.0 = 0.25$: a quarter of the neosome is
\omterm{def:b2:crustal-melting:migmatite}{leucosome}. (b) $0.25\times0.20 + 0.75\times4.0 = 0.05 + 3.00 = \qty{3.05}{\percent}$, as
in the mesosome: the closed-system interpretation holds. (c) For \qty{100}{\gram} of
parent, \qty{25}{\gram} of melt and \qty{75}{\gram} of residue formed; if $m$ grams of
melt left, $(25 - m)/(100 - m) = 0.10$, so $25 - m = 10 - 0.1m$ and $m = 15/0.9 =
\qty{16.7}{\gram}$: two thirds of the melt left the outcrop.
\end{solution}

\begin{solution}{pb:b2:crustal-melting:1}
\noindent\textbf{Part I.}

\textbf{1.} $39.10 + 3\times26.98 + 3\times28.09 + 12\times16.00 + 2\times1.008 =
\qty{398.33}{\gram\per\mole}$.

\textbf{2.} $18.02/398.33 = \qty{4.52}{\percent}$.

\textbf{3.} $200\times0.0452 = \qty{9.04}{\gram}$ of water per kilogram of schist.

\textbf{4.} $9.97 + (13.82 - 9.97)\times2/3 = \qty{12.54}{\percent}$. The least melt is
$9.04/0.1254 = \qty{72.1}{\gram}$ per kilogram: $F \ge \qty{7.21}{\percent}$
(\cref{prop:b2:crustal-melting:water}).

\textbf{5.} $9.04/0.060 = \qty{151}{\gram}$: $F = \qty{15.1}{\percent}$.

\textbf{6.} For example: the rock may hold less muscovite; part of the muscovite may
already have broken down below the solidus, with quartz, to K-feldspar and
aluminosilicate, its water leaving as a fluid before melting began; the melt may hold more
water than \qty{6}{\percent}; some water may stay in biotite.

\noindent\textbf{Part II.}

\textbf{7.} $680 - 7.5\times(7.8 - 2) = \qty{636.5}{\celsius}$; the rock peaks
\qty{43.5}{\celsius} below it and cannot melt, even with free water.

\textbf{8.} $605 + (19.2 - 15.3)/(27 - 15.3)\times65 = \qty{626.7}{\celsius}$. At
$t = 0$ the base, at \qty{595}{\celsius}, is \qty{32}{\celsius} below; at
\qty{10}{Myr}, at \qty{656}{\celsius}, it is above.

\textbf{9.} Melting at the wet solidus needs free water, and a heated metamorphic rock
holds very little: the free water is used up by the first melt, which then stops growing
until the micas break down.

\textbf{10.} Between 30 and \qty{40}{Myr}: $30 + (800 - 795)/(845 - 795)\times10 =
\qty{31.0}{Myr}$.

\textbf{11.} The base falls below \qty{12.5}{kbar} between 70 and \qty{80}{Myr}:
$70 + (13.7 - 12.5)/(13.7 - 12.4)\times10 = \qty{79.2}{Myr}$, at
$914 - 14\times0.923 = \qty{901}{\celsius}$, inside the amphibole bracket: an amphibolite
at the base would be melting.

\textbf{12.} $90 + (11.0 - 10)/(11.0 - 9.6)\times10 = \qty{97.1}{Myr}$, at $867 -
53\times0.714 = \qty{829}{\celsius}$. That is above the muscovite bracket: the muscovite
has broken down. It lies within the 800 to \qty{850}{\celsius} of biotite breakdown
(measured at lower pressure): biotite, the schist's other hydrous mineral, is at its
limit.

\noindent\textbf{Part III.}

\textbf{13.} $n(\ce{Al2O3}) = 15.0/101.96 = \qty{0.14712}{mol}$, $n(\ce{CaO}) =
0.55/56.08 = 0.00981$, $n(\ce{Na2O}) = 3.4/61.98 = 0.05486$, $n(\ce{K2O}) = 4.9/94.20 =
0.05202$.

\textbf{14.} A/CNK $= 0.14712/0.11669 = 1.261$.

\textbf{15.} A/NK $= 0.14712/0.10688 = 1.377$; A/CNK $> 1$: \omterm{def:b2:crustal-melting:asi}{peraluminous}.

\textbf{16.} \omterm{def:b2:crustal-melting:asi}{Peraluminous}, from a schist: an \omterm{def:b2:crustal-melting:types}{S-type granite}. Its excess aluminium sits in
muscovite and garnet (or cordierite).

\textbf{17.} $0.14712 - 0.11669 = \qty{0.03043}{mol}$ of \ce{Al2O3}; as muscovite,
$0.03043\times398.33 = \qty{12.1}{\gram}$ in \qty{100}{\gram} of granite.

\textbf{18.} The obsidian is (just) \omterm{def:b2:crustal-melting:asi}{metaluminous}, without the excess aluminium of a
metasedimentary melt: its source is probably different. But one sample proves little:
every granitoid suite has mildly \omterm{def:b2:crustal-melting:asi}{peraluminous} members, and a suite is judged by its whole
trend.

\noindent\textbf{Part IV.}

\textbf{19.} $10^{12}\,\unit{\cubic\metre}\times2650 = \qty{2.65e15}{kg}$.

\textbf{20.} The granite is the melt minus its water: $2.65\times10^{15}/0.94 =
\qty{2.82e15}{kg}$ of melt.

\textbf{21.} Only $0.151\times0.70$ of the schist's mass leaves as melt:
$2.82\times10^{15}/(0.151\times0.70) = \qty{2.67e16}{kg}$.

\textbf{22.} $2.67\times10^{16}/2800 = \qty{9.5e12}{\cubic\metre}$, that is
\qty{9.5e3}{\cubic\kilo\metre}, nearly ten times the \omterm{def:b2:magma-emplacement:bodies}{pluton}.

\textbf{23.} $9.5\times10^3/5 = \qty{1.9e3}{\square\kilo\metre}$, a square about
\qty{44}{km} on a side.

\textbf{24.} With $F = \qty{7.21}{\percent}$, the source is $15.1/7.21 = 2.1$ times
larger, about \qty{2.0e4}{\cubic\kilo\metre}. The volume of source crust melted to make
the \omterm{def:b2:magma-emplacement:bodies}{pluton}, with the \omterm{def:b2:mantle-melting:fraction}{melt fraction} of question~5, is
\textbf{\qty[text-series-to-math=true]{9.5e3}{\cubic\kilo\metre}}.
\end{solution}
